% Put the project assessment here

\section{Project Assessment}

\subsection{User Testing}

User testing was conducted to verify the functional requirements of the Digital Farmers' Produce System, assess its non-functional qualities, identify defects, and determine whether the intended users could use the system effectively. The testing involved representative users from the three main roles of the system: farmers, buyers, and Department of Agriculture/Municipal Agriculture Office (DA/MAO) staff.

The pilot assessment included twelve (12) respondents composed of five (5) farmers, four (4) buyers, and three (3) DA/MAO officers. These users were selected because they represent the actual stakeholders who will interact with the system during deployment.

\subsubsection{Testing Objectives}

The user testing was conducted with the following objectives:

\begin{enumerate}
    \item To verify that each functional requirement of the Digital Farmers' Produce System works as intended.
    \item To evaluate non-functional requirements such as usability, performance, reliability, compatibility, and responsiveness.
    \item To detect bugs, usability issues, and workflow gaps before deployment.
    \item To determine whether the system is acceptable and meaningful to its intended users.
\end{enumerate}

The overall aim was to ensure that the system was usable, reliable, and acceptable to the users, and that no critical defects remained in the core workflows prior to implementation.

\subsubsection{Target Users and Test Environment}

Testing was conducted using devices that approximate the expected real-world environment of the target users. Low-end Android smartphones were used to represent the typical devices of farmers and buyers, while a standard desktop or laptop browser was used to reflect the working environment of DA/MAO staff.

The system was tested using a staging version of the Digital Farmers' Produce System with sample test data. This was done to protect real user information and agricultural records while still allowing users to perform realistic tasks such as registration, product posting, searching, messaging, request submission, and dashboard monitoring.

\subsubsection{Functional Test Specifications}

Functional testing was scenario-based. Each test case included a test ID, description, preconditions, step-by-step procedure, expected result, actual result, and pass/fail status. The test scenarios were grouped according to the major user roles of the system.

\begin{table}[H]
\centering
\caption{Functional Test Scenarios for Farmers}
\begin{tabular}{p{4cm}p{6cm}p{3cm}}
\hline
\textbf{Test Scenario} & \textbf{Expected Result} & \textbf{Result} \\
\hline
Register a farmer account & The farmer account is successfully created after valid information is submitted. & Passed \\
Log in to the system & The farmer is redirected to the farmer dashboard. & Passed \\
Create a crop listing & The crop listing is saved and displayed in the farmer's product list. & Passed \\
Edit an existing listing & The updated crop details are reflected in the system. & Passed \\
Deactivate a listing & The selected listing is removed from active listings. & Passed \\
View and reply to messages & The farmer can open and respond to buyer messages. & Passed \\
Submit a resource request & The request is recorded and becomes visible to DA/MAO staff. & Passed \\
\hline
\end{tabular}
\end{table}

\begin{table}[H]
\centering
\caption{Functional Test Scenarios for Buyers}
\begin{tabular}{p{4cm}p{6cm}p{3cm}}
\hline
\textbf{Test Scenario} & \textbf{Expected Result} & \textbf{Result} \\
\hline
Register a buyer account & The buyer account is successfully created after valid information is submitted. & Passed \\
Log in to the system & The buyer is redirected to the buyer dashboard. & Passed \\
Search for products & Relevant product listings are displayed based on the search term. & Passed \\
Filter product listings & The system displays products based on the selected category, location, or availability. & Passed \\
Open a product listing & The buyer can view product details and farmer information. & Passed \\
Send a message to a farmer & The message is successfully sent to the selected farmer. & Passed \\
\hline
\end{tabular}
\end{table}

\begin{table}[H]
\centering
\caption{Functional Test Scenarios for DA/MAO Staff}
\begin{tabular}{p{4cm}p{6cm}p{3cm}}
\hline
\textbf{Test Scenario} & \textbf{Expected Result} & \textbf{Result} \\
\hline
Log in as DA/MAO staff & The DA/MAO dashboard is displayed. & Passed \\
View registered farmers & The list of registered farmers is displayed. & Passed \\
View registered buyers & The list of registered buyers is displayed. & Passed \\
Review resource requests & Submitted farmer requests are displayed for review. & Passed \\
Approve or reject requests & The request status is updated successfully. & Passed \\
Export records to CSV & The system generates a downloadable CSV file. & Passed \\
\hline
\end{tabular}
\end{table}

A functional requirement was considered passed when the task could be completed without system error, the actual result matched the expected result, and the assigned users were able to complete the scenario with minimal assistance.

\subsubsection{Non-Functional Testing}

Non-functional testing was conducted to evaluate usability, performance, reliability, compatibility, and responsiveness.

\begin{table}[H]
\centering
\caption{Non-Functional Testing Criteria}
\begin{tabular}{p{4cm}p{7cm}p{3cm}}
\hline
\textbf{Criteria} & \textbf{Description} & \textbf{Result} \\
\hline
Usability & Users were observed while completing tasks such as posting, searching, messaging, and submitting requests. & Satisfactory \\
Performance & Page loading and response times were observed during login, dashboard loading, posting, searching, and messaging. & Satisfactory \\
Reliability & The system was monitored for crashes, unhandled errors, and repeated failures. & Satisfactory \\
Compatibility & The system layout was checked on Android smartphones, desktops, and laptop browsers. & Satisfactory \\
Responsiveness & The interface was checked for proper display on different screen sizes. & Satisfactory \\
\hline
\end{tabular}
\end{table}

The results showed that the system remained usable and stable during testing. Users were able to complete the core workflows, and no critical crashes were observed. The interface adjusted properly on both mobile and desktop screens.

\subsubsection{Bug Detection and Tracking}

Bugs and usability issues were identified through developer testing, observation during user sessions, and feedback from the respondents. A defect log was used to record all issues found during testing. Each defect entry included the following details:

\begin{itemize}
    \item Unique defect ID
    \item Date reported
    \item Module or page affected
    \item Severity level
    \item Steps to reproduce
    \item Expected result
    \item Actual result
    \item Assigned developer
    \item Status
\end{itemize}

Issues were classified as critical, major, or minor. Critical issues refer to defects that prevent users from completing a core function. Major issues affect important workflows but have possible workarounds. Minor issues refer to layout, labeling, or usability concerns that do not prevent task completion.

\subsubsection{Bug Fixing and Regression Testing}

After the bugs and usability issues were recorded, the identified problems were prioritized and corrected. Examples of fixes included improving unclear labels, adjusting layout elements for mobile screens, validating required fields, and refining error messages.

Regression testing was conducted after the fixes were implemented. The affected test cases were repeated to verify that the corrections worked properly. Core workflows such as login, product posting, searching, messaging, resource request submission, and dashboard monitoring were also retested to ensure that no new errors were introduced.

Before deployment, the target condition was to have no remaining critical defects and no remaining major defects in the core workflows. Minor issues, if any, were documented for future enhancement.

\subsubsection{User Acceptance Testing}

User Acceptance Testing was conducted using the Technology Acceptance Model (TAM). The evaluation focused on the following constructs:

\begin{itemize}
    \item Perceived Usefulness
    \item Perceived Ease of Use
    \item Attitude Toward Using
    \item Behavioral Intention to Use
\end{itemize}

Additional survey items were also included to assess system functionality and overall satisfaction. The respondents answered a five-point Likert-scale questionnaire immediately after hands-on testing. The scale ranged from 1, interpreted as Strongly Disagree, to 5, interpreted as Strongly Agree.

The summarized results were analyzed using means, standard deviations, and descriptive interpretation. High mean values indicate strong user acceptance, while lower values indicate areas that may require improvement. Since only aggregated results were supplied for the manuscript, internal consistency reliability using Cronbach's alpha was not calculated.

\subsection{Security Testing}

Security testing focused on ensuring the confidentiality, integrity, and availability of user data and system functions. The security assessment was designed to identify possible vulnerabilities before full deployment.

The security testing process included configuration review, code review, automated vulnerability scanning, and manual security checks. The system was also reviewed based on common web application security risks.

\subsubsection{Security Testing Tool}

The system will be tested using OWASP Zed Attack Proxy (ZAP), an open-source web application security scanner. OWASP ZAP is used to identify common vulnerabilities such as SQL injection, cross-site scripting, insecure cookies, weak session handling, and missing security headers. The tool can be accessed at \href{https://www.zaproxy.org/}{https://www.zaproxy.org/}.

\subsubsection{Security Test Coverage}

The security testing covered the following areas:

\begin{itemize}
    \item Verification that the system uses SHA 256 for encrypted communication.
    \item Review of password storage and authentication mechanisms.
    \item Checking of server-side role-based access control.
    \item Testing of form inputs against SQL injection attempts.
    \item Testing of text fields against cross-site scripting attempts.
    \item Checking of session handling and logout behavior.
    \item Review of error messages to ensure that internal system details are not exposed.
    \item Checking for missing or weak security headers.
\end{itemize}

\subsubsection{Possible Security Vulnerabilities}

The following security vulnerabilities were considered during assessment:

\begin{table}[H]
\centering
\caption{Security Vulnerabilities and Mitigation Measures}
\begin{tabular}{p{4cm}p{5cm}p{5cm}}
\hline
\textbf{Possible Vulnerability} & \textbf{Risk} & \textbf{Mitigation Measure} \\
\hline
SQL Injection & Attackers may attempt to manipulate database queries through input fields. & Use prepared statements, input validation, and database query parameterization. \\
Cross-Site Scripting (XSS) & Attackers may attempt to inject malicious scripts into forms or messages. & Apply output escaping, input sanitization, and content security policies. \\
Weak Authentication & Unauthorized users may attempt to access protected accounts. & Enforce password hashing, strong password rules, and secure login validation. \\
Broken Access Control & Users may attempt to access pages or records outside their role. & Apply server-side role checks for every protected route and resource. \\
Insecure Session Handling & Sessions may be reused or hijacked if not managed properly. & Use secure session configuration, logout invalidation, and HTTPS-only cookies. \\
Missing Security Headers & Browsers may not receive instructions to prevent common attacks. & Configure headers such as Content-Security-Policy, X-Frame-Options, and X-Content-Type-Options. \\
Information Disclosure & Error messages may reveal database details, file paths, or internal configuration. & Replace technical errors with generic user-friendly messages and log detailed errors on the server. \\
\hline
\end{tabular}
\end{table}

\subsubsection{Manual Security Checks}

Manual testing was also included to supplement the automated scan. The following checks were performed or planned:

\begin{itemize}
    \item Attempting to modify URL parameters to access another user's record.
    \item Entering script tags in input fields to check for XSS vulnerabilities.
    \item Submitting invalid form data to verify validation and error handling.
    \item Attempting to access protected pages without logging in.
    \item Attempting to access role-specific pages using a different user role.
\end{itemize}

A security issue is considered present if the scanner detects medium- or high-severity vulnerabilities, if unauthorized access becomes possible, if inputs are interpreted as executable content, or if error messages expose internal database or server details.

\subsubsection{Mitigation and Re-Testing}

When vulnerabilities are found, the following mitigation actions will be applied:

\begin{itemize}
    \item Strengthen input validation and output escaping.
    \item Use prepared statements for all database transactions.
    \item Improve password policies and password hashing.
    \item Add CSRF tokens to forms that modify system data.
    \item Enforce server-side role-based access control.
    \item Configure recommended HTTP security headers.
    \item Replace detailed technical errors with generic user-facing messages.
\end{itemize}

After applying the fixes, automated scanning using OWASP ZAP and manual testing will be repeated. This regression security testing will confirm whether the vulnerabilities were removed and whether the fixes introduced any new issues.

The security testing phase will be considered satisfactory when no high-severity vulnerabilities remain, manual attempts to bypass access control fail, and the system behaves securely and predictably under the tested conditions.